% p1-02-background

\section{P1 §2. Background and Related Work}

2. Background and Related Work

       Figure (fig-p1-ch2-background). Background triptych --- History Scroll / MDL Scales / Trinity
                                               Anchor.

  2.1 Symbolic Approaches to Physical Constants: Historical Context
  Attempts to express dimensionless physical constants in closed algebraic form have a
  long and largely undisciplined history. Early efforts concentrated on alpha-1 approx
  137, with various proposals involving integers and simple arithmetic --- all post-hoc and
  statistically uncontrolled (PDG 2024 review articles). The combinatorial proliferation of
  candidate expressions at increasing complexity was universally ignored, making
  significance claims meaningless.

  More systematic approaches employing continued fractions, Ramanujan-type
  hypergeometric identities, and modular forms have examined the representability of

  specific mathematical constants such as pi and e. These efforts differ from the present
  work in that they target mathematical identities rather than empirical physical
  constants, and they do not formulate the problem as a statistical model-selection
  problem over a structured hypothesis class with explicit null models. The present paper
  belongs to a different methodological tradition.

  2.2 Symbolic Regression in Physics
  Symbolic regression --- the automated discovery of mathematical expressions fitting
  numerical data --- has been developed extensively in the machine learning community,
  with methods ranging from genetic programming to neural-guided search. The AI
  Feynman algorithm of Udrescu and Tegmark (Science Advances, 2020) demonstrated
  that physics-inspired techniques including dimensional analysis, separability detection,
  and brute-force enumeration can recover 100 equations from the Feynman Lectures on
  Physics. More recent work in physics-informed symbolic regression (Angelis et al.,
  Archives of Computational Methods in Engineering, 2023) has extended these methods
  to broader scientific domains.

  The present work differs from all these approaches in a fundamental respect: the target
  is not a functional relationship between variables but a set of isolated numerical values
  of dimensionless constants, and the search is not over a continuous parameter space
  but over a discretely bounded combinatorial grammar. This distinction is crucial for the
  statistical methodology: the absence of variables means that null models can be
  precisely specified and that the effective trial count is the cardinality of the discrete
  hypothesis class.

  2.3 Minimum Description Length and Bayesian Model Comparison
  The Minimum Description Length principle, originating in Rissanen's 1978 paper
  Modeling by the shortest data description (Automatica, 1978) and systematized in
  Barron, Rissanen \& Yu (1998), provides a rigorous basis for model comparison that
  penalizes complexity in terms of code length. The core insight is that any regularity in
  data can be exploited to compress it, and the best model is the one achieving the
  shortest description of data and model combined. Grünwald's comprehensive treatment
  (MIT Press, 2007) formalizes the connections between crude two-part codes,
  normalized maximum likelihood, Bayesian marginal likelihoods, and prequential coding,
  establishing asymptotic equivalences between these approaches.

  Bayesian model comparison via marginal likelihoods provides an equivalent framework
  when proper priors are specified; the resulting Bayes factor quantifies the relative
  evidence in favor of structured over unstructured models. Both frameworks are
  deployed in Section 7 with explicit penalization for the size and structure of the
  symbolic hypothesis class.

  The connection to Kolmogorov-Chaitin algorithmic complexity (Chaitin, Cambridge
  University Press, 1987) is conceptual: the MDL compression score DeltaMDL provides a
  computable, statistical lower bound on the potential reduction in Kolmogorov
  complexity achievable by the structured grammar.

  2.4 The Algebraic Anchor: Lucas Ring and Trinity Identity

  The choice of $\varphi$ as a basis generator is anchored in the algebraic identity

  $\varphi$2 + $\varphi$-2 = 3.

  which holds exactly because $\varphi$ satisfies the minimal polynomial $\varphi$2 = $\varphi$ + 1. This
  identity --- the Trinity identity of the broader research program --- establishes that the
  squared golden ratio and its reciprocal square sum to the integer 3, implying a ternary
  cardinality relation within the Lucas ring Z[$\varphi$]. In the context of the present paper, its
  significance is limited to the following algebraic observation: Z[$\varphi$] is closed under
  multiplication, ensuring that integer-exponent products of $\varphi$ remain in a
  well-characterized algebraic structure. This algebraic closure simplifies the linear
  independence argument for the log-lattice (Section 3.3) and guarantees that the
  hypothesis class HC is non-degenerate. The identity is provenance from the Trinity $S^{3}$AI
  PhD monograph (PhD: Ch.3 Trinity Identity); no further physical significance is claimed
  for it in this paper.

  2.5 Positioning
  The present work occupies a specific methodological niche: rigorous combinatorial and
  information-theoretic evaluation of whether a finite, pre-specified symbolic language
  achieves statistically non-random compressibility over a frozen, pre-committed dataset
  of dimensionless constants. The novelty lies in (i) the explicit treatment of the search as
  a multiple-testing problem, (ii) the use of MDL with a complexity-penalized grammar
  model, (iii) the Pellis hierarchical expansion as a two-tier refinement mechanism, (iv)
  the provision of three complementary null models and five explicit falsification criteria,
  and (v) the transparent provenance connection to the PhD monograph, with clear
  delineation of what the monograph contributes and what it does not validate.

  2.6 Reviewer-Hardening Protocol Imported from the PhD Monograph
  The Flos Aureus PhD monograph contributes a procedural discipline that is useful for
  this paper but does not count as evidence for the paper's empirical claims. Four rules
  are adopted:

  1. No silent overclaim. A symbolic fit is never promoted to a physical derivation unless a
  mechanism is supplied.
  2. No stale constants. Public constants are locked to the cited PDG/CODATA release,
  and any updated value is treated as an out-of-sample update rather than retrofitted into
  the original claim.
  3. No hidden search space. Every reported fit is paired with an explicit hypothesis-class
  cardinality Neff=|S(C)|, so the look-elsewhere correction is visible.
  4. No proof-status inflation. Coq-backed algebraic identities may be cited as formal
  provenance only at their actual status; Admitted material is treated as a proof sketch,
  not as a theorem.

  This protocol is intentionally conservative. Its purpose is to make the article easier to
  criticize correctly and harder to misread as numerology, proof-by-catalogue, or a claim
  that the Standard Model constants have been derived from $\varphi$.

  2.7 Claim-Tier Tags Imported from the PhD Defense Protocol
  To prevent category errors, every claim in the article is interpreted through one of four
  tiers:

  - Tier A --- algebraic identity: exact symbolic statements such as $\varphi$2+$\varphi$-2=3. These are
  mathematical claims.
  - Tier B --- symbolic compressibility: low-description-length approximations under a
  fixed grammar and a fixed dataset. These are statistical claims and require null models.
  - Tier C --- mechanism candidate: a proposed bridge from algebra to physics. These are
  conjectural unless embedded in an accepted physical theory.
  - Tier D --- historical or interpretive context: conceptual framing, including
  golden-section history and visual architecture. These passages provide orientation only
  and are not evidence for Tier B or Tier C.

  This tier system is deliberately reviewer-facing. If a sentence can be read in more than
  one tier, the lower evidentiary tier controls the interpretation.

\subsection{References}

- Vasilev, D. \textit{Flos Aureus Monograph: Trinity Constants}. DOI \href{https://zenodo.org/doi/10.5281/zenodo.19227877}{10.5281/zenodo.19227877}.
- Pellis, S. "The Fine-Structure Constant and the Golden Ratio." \href{https://vixra.org/pdf/2110.0117v4.pdf}{viXra 2110.0117 v4}.
- Sherbon, M. A. "Fundamental Nature of the Fine-Structure Constant." \href{https://hal.science/hal-02341850/document}{HAL hal-02341850}.
